\begin{claim}{F.21}{Proved}{Full cross-divisor transitivity and normal generation inside the joint kernel}
Retain the actual groups, source maps, eight selected positive
cross-collision meridians, and base whiskers of F.17--F.19.
Write
\[
 G=P_{10}^{\dagger},\qquad
 \varpi:G\twoheadrightarrow\mathcal Q,
 \qquad D^+=\ker\varpi,\qquad
 H^\dagger=\Phi^\dagger(L^\dagger),\qquad
 K^+=H^\dagger\cap D^+ .
                                                               \tag{F.21.1}
\]
For each $u\in\mu _5\setminus\{\zeta\}$ and
$\epsilon\in\{1,-1\}$, let $\mathcal P_{u,\epsilon}$
be the full closed image in $\mathcal Q$ of its actual
cross divisor $v_u=\epsilon a$, based using the source
whisker selected for $\nu_{u,\epsilon}$ in F.19.
Then the source subgroup $\gamma(L)$ has a transitive
left action on the entire compact coset space
\[
                 \mathcal Q/\mathcal P_{u,\epsilon}.
                                                               \tag{F.21.2}
\]
This holds for each of the eight named divisors, with
the original label maps and full profinite completions.
The corresponding discrete coset action is also transitive.

The selected meridian has the following actual
centralizer-image containment:
\[
 \varpi\bigl(C_G(\nu_{u,\epsilon})\bigr)
       \supseteq\mathcal P_{u,\epsilon}.
                                                               \tag{F.21.3}
\]
Write $D^+_{\mathrm{nil}}$ for the inverse limit of all
continuous finite nilpotent quotients of $D^+$.
Consequently
\[
 \begin{gathered}
 \overline{\langle\!\langle K^+\rangle\!\rangle}_{D^+}=D^+,\\
 K^+\twoheadrightarrow
       D^+_{\mathrm{ab}}
       :=D^+/\overline{[D^+,D^+]},\qquad
 K^+\twoheadrightarrow D^+_{\mathrm{nil}} .
 \end{gathered}
                                                               \tag{F.21.4}
\]
In particular every continuous finite nilpotent quotient of
$D^+$ has the same image on $K^+$ and on $D^+$.
The quotient maps and the normal closure are the displayed
actual maps and subgroups.
The remaining equality $K^+=D^+$ and the source-membership
equation of F.19 remain unresolved.
\end{claim}

\begin{proof}
\emph{Actual source basepoints on each filled cross divisor.}
Put
\[
 \begin{gathered}
 A(s)=\frac{\zeta-s}{\zeta+s},\qquad
 C_u(s)=\frac{u-s}{u+s},\qquad B=A^2,\qquad U=C_u^2,\qquad
 V=\epsilon A,\\
 k_u=\frac{u-\zeta}{u+\zeta},\qquad
 \mathscr A=\mathbb P^1\setminus\{0,\infty,1,-1\}.
 \end{gathered}                                                \tag{F.21.5}
\]
For the selected meridian $\nu_{u,\epsilon}$, let $d$ be
its collision point in $Y$. Thus $d^2=A_u$ for
$\epsilon=1$, and $d^2=B_u$ for $\epsilon=-1$, in the
notation of F.17.2. The sixteen collision points are
distinct and lie in $Y$ by F.17. At $d$ exactly this
cross collision occurs. The filled map
$\gamma:Y\to X$ of F.19.15 is defined at $d$; its
selected two-track projection is
\[
       s\longmapsto\bigl(A(s)^2,C_u(s)^2,\epsilon A(s)\bigr).
                                                               \tag{F.21.6}
\]
Its value at $d$ is on the diagonal $U=V$.

Keep the actual source whisker $\omega$ of that meridian,
extended to its endpoint $d$. It is a path in $Y$ from
the original $s_0$ to $d$, with its open terminal part
in $Y^\dagger$. The based source isomorphism is
$g\mapsto\omega g\omega^{-1}$ from $\pi_1(Y,d)$ to
$\pi_1(Y,s_0)$. On $X$ it is conjugation by the literal
path $\gamma(\omega)$. All forgetting projections
commute with these maps, because they forget coordinates
of this same configuration path. We prove transitivity
at the basepoint $d$ and then apply these specified
isomorphisms. The cross-divisor subgroup so obtained is
exactly the one based by the chosen whisker. No change
of the selected meridian or its conjugacy coordinate
is required.

\emph{The fourfold square-root cover and its pencil.}
The common two-coordinate moduli space is
\[
 \mathcal M_{0,5}
       =\{(B,U): B\notin\{0,1\},\ U\notin\{0,1,B\}\}.
\]
Both coordinates are finite. Its actual connected
fourfold covering and an intermediate double covering are
\[
 \begin{gathered}
 W=\{(A,C)\in\mathscr A^2:C\ne A,\ C\ne-A\},\\
 W_A=\{(A,U):A\in\mathscr A,\ U\notin\{0,1,A^2\}\},\\
 W\xrightarrow{\rho}W_A\xrightarrow{\beta}\mathcal M_{0,5},\\
 \rho(A,C)=(A,C^2),\qquad\beta(A,U)=(A^2,U).
 \end{gathered}                                               \tag{F.21.7}
\]
These maps are coverings because the squared coordinates
avoid $0$ and $\infty$; their four, respectively two,
local square-root choices are distinct. The complements
are connected: their projections to $A\in\mathscr A$
have connected punctured-plane fibers, and local
trivializations arise by moving the finitely many
excluded points in disjoint small disks.

There is an everywhere-defined map
\[
 \kappa:W\longrightarrow
        \mathscr B:=\mathbb P^1\setminus\{0,1,-1\},
       \qquad \kappa(A,C)=\frac{C-A}{1-AC}.
                                                               \tag{F.21.8}
\]
The numerator and denominator cannot vanish together on
$W$: that would give $A=C$ and $A^2=1$.
The values $0,1,-1$ would force, respectively,
$C=A$, $(C-1)(A+1)=0$, or $(C+1)(1-A)=0$.
All are excluded in (F.21.7).
Let
\[
 E=\{AC=1\}\subset W,
 \qquad W^\circ=W\setminus E,
 \qquad\mathscr B^\circ=\mathscr B\setminus\{\infty\}.
                                                               \tag{F.21.9}
\]
For finite $k\in\mathscr B^\circ$, the inverse coordinate
formula is
\[
 C=\frac{A+k}{1+kA},\qquad
 A\in\mathbb P^1\setminus
 \{0,\infty,1,-1,-k,-1/k,r_+(k),r_-(k)\},
 \quad kr^2+2r+k=0.
                                                               \tag{F.21.10}
\]
The eight displayed punctures are distinct for
$k\notin\{0,1,-1,\infty\}$.
Indeed the two roots are distinct by $1-k^2\ne0$;
neither is $0$ or $\infty$; the substitutions $r=1,-1$
give $2(k+1),2(k-1)$; and the substitutions
$r=-k,-1/k$ give $k(k^2-1),(k^2-1)/k$.
The remaining six punctures are pairwise distinct
under the same exclusions.

On a small simply connected neighborhood of any such
$k$, choose the two distinct analytic roots $r_+,r_-$.
All eight puncture paths remain in pairwise disjoint
coordinate disks. Within each disk, translate its
puncture using a bump function that is one near the
path and zero near the disk boundary. These isotopies,
with the inverse isotopies, depend continuously on $k$.
Their disjoint supports produce an actual local
trivialization of the complements in (F.21.10).
Consequently $W^\circ\to\mathscr B^\circ$ is a locally
trivial bundle with the indicated eight-punctured sphere
fiber. Its base and fiber are connected, and its
fundamental-group sequence is exact.

The old source has constant pencil value:
\[
 \kappa\bigl(A(s),C_u(s)\bigr)=k_u.
                                                               \tag{F.21.11}
\]
Here $k_u\notin\{0,1,-1,\infty\}$ since $u\ne\zeta$,
$u\ne-\zeta$, and both $u,\zeta$ are nonzero.
The coordinate $s=\zeta(1-A)/(1+A)$ identifies its
entire fiber with
\[
 Y_u=\mathbb P^1\setminus
       \{0,\infty,\zeta,-\zeta,u,-u,t_u,-t_u\},
       \qquad t_u^2=\zeta u .
                                                               \tag{F.21.12}
\]
For $u=\zeta^j$, take $t_u=\zeta^h$ with
$h=3(1+j)\pmod5$. The four exponent pairs are
\[
                 (j,h)\in\{(0,3),(2,4),(3,2),(4,0)\}.
\]
The inverse coordinate in (F.21.10) becomes exactly
$C_u(s)$. Its exclusions are $s=0,\infty,\pm\zeta$
from $A\in\mathscr A$, $s=\pm u$ from $C_u\in\mathscr A$,
and $s=\pm t_u$ from $C_u=-A$; $C_u=A$ adds only
$0,\infty$. Thus $Y\to Y_u$ fills exactly four of
the old punctures, and induces a surjection on
fundamental groups. Every loop on $Y_u$ can be moved
away from those four points by a small path homotopy.
The actual image of $\pi_1(Y,d)$ in $\pi_1(W^\circ)$
is therefore the image of this whole bundle fiber.

\emph{Filling the connected multiplicity-one divisor.}
The divisor $E$ is the actual connected smooth curve
\[
 A\longmapsto(A,A^{-1}),\qquad
 A\in\mathbb P^1\setminus\{0,\infty,1,-1,i,-i\}.
                                                               \tag{F.21.13}
\]
The extra exclusions $\pm i$ express $C\ne-A$.
At $E$ the local coordinate at the target infinity is
\[
 \frac1\kappa=\frac{1-AC}{C-A},\qquad
 \left.\frac{\partial(1/\kappa)}{\partial C}\right|_{C=1/A}
       =\frac{-A}{A^{-1}-A}\ne0.
                                                               \tag{F.21.14}
\]
Thus a positive transverse meridian of $E$ maps to
a single positive meridian of infinity in the base.

Set $J^\circ=\pi_1(W^\circ,(A(d),C_u(d)))$ and
$J=\pi_1(W,(A(d),C_u(d)))$. The inclusion is onto
on fundamental groups, and its kernel is the normal
closure of the meridians of $E$. To see both statements
directly, perturb a loop in the smooth manifold $W$
away from the real-codimension-two submanifold $E$.
Perturb a null-homotopy disk relative to its boundary
to meet $E$ transversely in finitely many interior
points. Delete small disks around those points;
the resulting punctured disk expresses the original
loop as a product of conjugates of the transverse
meridians, with their intersection signs.
All positive meridians form one conjugacy class,
because $E$ is connected and its complex normal line
has its fixed orientation.

Let $N^\circ$ be the fiber image in $J^\circ$ and let
$N$ be its image in $J$. The bundle sequence and
(F.21.14) give the following quotient map with its
specified inverse identification:
\[
 \begin{aligned}
 J/N
 &\cong
 (J^\circ/N^\circ)/
       \langle\!\langle\hbox{meridian of }\infty
                                \rangle\!\rangle\\
 &\cong
 \pi_1(\mathscr B^\circ,k_u)/
       \langle\!\langle\hbox{meridian of }\infty
                                \rangle\!\rangle
 \cong\pi_1(\mathscr B,k_u).
 \end{aligned}                                                \tag{F.21.15}
\]
The maps are induced by $\kappa$ and the two actual
fillings. Their kernels are the normal subgroups
just specified, so the inverse maps are the unique
factorizations of these same quotient homomorphisms.
In particular
\[
 N=\operatorname{im}\bigl(\pi_1(Y,d)\to J\bigr)
      =\ker\bigl(\kappa_*:J\twoheadrightarrow
                                  \pi_1(\mathscr B,k_u)\bigr).
                                                               \tag{F.21.16}
\]
This is an equality in the discrete fundamental group
of the named covering. No assertion about completing
this kernel sequence is needed below.

\emph{The new section and the onto pencil map.}
Base $W_A$ at $(A(d),C_u(d)^2)=(A(d),\epsilon A(d))$
and put $T=\pi_1(W_A)$.
The cover $W\to W_A$ embeds $J$ as a normal subgroup
of index two. The map $(A,U)\mapsto A$ has the actual
based section
\[
       \sigma_\epsilon:\mathscr A\longrightarrow W_A,
                    \qquad A\longmapsto(A,\epsilon A).
                                                               \tag{F.21.17}
\]
Its image $M=\sigma_{\epsilon,*}\pi_1(\mathscr A,A(d))$
maps isomorphically to $\pi_1(\mathscr A,A(d))$ under
the projection. The kernel of that projection is
the actual one-point-pushing group
\[
 F=\pi_1\bigl(\mathbb P^1\setminus
                   \{0,\infty,1,A(d)^2\},\epsilon A(d)\bigr)
       \cong F_3 .
                                                               \tag{F.21.18}
\]
Its injection and kernel description follow from the
configuration bundle, because $\pi_2(\mathscr A)=0$.

The pullback of $W\to W_A$ by this very section is
the connected double covering
\[
 \mathscr R=\mathbb P^1\setminus\{0,\infty,1,-1,i,-i\}
       \longrightarrow\mathscr A,
              \qquad r\longmapsto\epsilon r^2,
                                                               \tag{F.21.19}
\]
and its map into $W$ is $r\mapsto(\epsilon r^2,r)$.
Its basepoint is exactly $r_0=C_u(d)$; the collision
equation gives $A(d)=\epsilon r_0^2$.
Thus its fundamental-group image is exactly $M\cap J$.
It also proves that the projection of $M$ onto $T/J$
is onto. On this pullback the pencil map is literally
\[
 f_\epsilon(r)=\kappa(\epsilon r^2,r)
               =\frac{r}{r^2+\epsilon r+1}.
                                                               \tag{F.21.20}
\]
At $r_0$ it has value $k_u$, with no intervening
basepoint path.

This map induces a surjection
$\pi_1(\mathscr R,r_0)\to\pi_1(\mathscr B,k_u)$.
Here is a direct covering proof. Its derivative is
$(1-r^2)/(r^2+\epsilon r+1)^2$, and the exact identities
\[
 \begin{aligned}
 f_\epsilon(r)-\epsilon/3
     &=\frac{-\epsilon(r-\epsilon)^2}
                       {3(r^2+\epsilon r+1)},\\
 f_\epsilon(r)+\epsilon
     &=\frac{\epsilon(r+\epsilon)^2}{r^2+\epsilon r+1},\\
 f_\epsilon(r)-\epsilon
     &=\frac{-\epsilon(r^2+1)}{r^2+\epsilon r+1}
 \end{aligned}                                                \tag{F.21.21}
\]
identify every removed preimage and both branch values.
In particular
\[
 f_\epsilon:\mathscr R\longrightarrow
       \mathscr B\setminus\{\epsilon/3\}
                                                               \tag{F.21.22}
\]
is a connected unramified degree-two covering.
The two preimages of $0$ are $0,\infty$; the preimages
of $\epsilon$ are $i,-i$; the single double preimage
of $-\epsilon$ is $-\epsilon$; and the single double
preimage of $\epsilon/3$ is $\epsilon$.
The poles of $f_\epsilon$ map to the allowed target
infinity and are ordinary unramified points.
These are precisely the six punctures of $\mathscr R$.

Write $G_0=\pi_1(\mathscr B\setminus\{\epsilon/3\},k_u)$
and let $\chi:G_0\to C_2$ be the monodromy character
of (F.21.22). Its index-two kernel is the image of
$\pi_1(\mathscr R,r_0)$.
A based meridian $c$ of $\epsilon/3$ has $\chi(c)=1$
by its local double-root expression in (F.21.21).
The filling $q:G_0\twoheadrightarrow\pi_1(\mathscr B,k_u)$
kills $c$. For every $g\in G_0$, the element
$g c^{-\chi(g)}$ belongs to $\ker\chi$ and has the same
image under $q$ as $g$. Hence $q(\ker\chi)$ is the
entire target. This proves the asserted surjectivity,
using the actual index-two cover and actual filling.
The basepoint avoids $\epsilon/3$: its fiber in
(F.21.21) consists only of $r=\epsilon$, which is
excluded from $\mathscr R$.

Combining this surjectivity with (F.21.16) proves the
exact subgroup factorizations
\[
                       J=N(M\cap J),\qquad T=JM=NM .
                                                               \tag{F.21.23}
\]
For the first, given $j\in J$, choose $m\in M\cap J$
with $\kappa_*(m)=\kappa_*(j)$ and write
$j=(jm^{-1})m$ with $jm^{-1}\in N$.
For the second, $M\to T/J$ is onto by the connected
cover (F.21.19), and $(M\cap J)M=M$.

\emph{The exact orbit map.}
At the diagonal basepoint $d$, the two copies of
$\pi_1(\mathcal M_{0,5})$ have exactly the same
base configuration. Let $P_{5,\mathrm d}$ and
$P_{4,\mathrm d}$ denote these discrete groups with
their actual forgetting homomorphism. The mutually
inverse diagonal-coset maps are the following, where
$[b_o,b_n]$ denotes the indicated right diagonal coset:
\[
 \begin{gathered}
 \mathcal Q_{\mathrm d}
     =P_{5,\mathrm d}\mathbin{\times}_{P_{4,\mathrm d}}P_{5,\mathrm d},
 \qquad
 \mathcal Z_{\mathrm d}
     =\mathcal Q_{\mathrm d}/\operatorname{diag}(P_{5,\mathrm d}),\\
 \mathcal Z_{\mathrm d}\longrightarrow F,
        \qquad[b_o,b_n]\longmapsto b_nb_o^{-1},\\
 F\longrightarrow\mathcal Z_{\mathrm d},
        \qquad v\longmapsto[1,v].
 \end{gathered}                                               \tag{F.21.24}
\]
Their composites are the identity, by right
multiplication with $(b_o^{-1},b_o^{-1})$.
The group $F$ here is the same one-point kernel
(F.21.18); passage to the cover $W_A$ does not
alter its fiber or its inclusion.

For $n\in N$, write $a(n)\in\pi_1(\mathscr A)$ for
its actual $A$-coordinate and let $m(a(n))\in M$
be its image under (F.21.17).
The actual source pair has precisely the image
\[
                   n\longmapsto(n,m(a(n)))\qquad(n\in N).
                                                               \tag{F.21.25}
\]
Indeed its old projection has image $N$ by
(F.21.16), and its new projection is
$s\mapsto(A(s)^2,\epsilon A(s))$, the displayed
section applied to the same $A$-loop.
Any element killed by the old projection is also
killed by its $A$-projection and hence by the new
projection. This proves the equality of the source
image with (F.21.25), including its maps.

Left multiplication on (F.21.24) acts on $F$ as
\[
                      v\longmapsto m(a(n))\,v\,n^{-1}.
                                                               \tag{F.21.26}
\]
For every $v\in F$, factor $v^{-1}=nm$ using
(F.21.23). Projection to $\pi_1(\mathscr A)$ gives
$a(m)=a(n)^{-1}$. The projection is an isomorphism
on $M$, so $m=m(a(n))^{-1}$ and
$v=m(a(n))n^{-1}$. This is the image of $1$ under
(F.21.26). The source action is therefore transitive
on the entire discrete coset space.

The connected-fiber forgetting map from the F.19 space
$X$ to this selected two-track space carries the
selected divisor onto its diagonal. Its divisor
image is the full inverse image of the diagonal
group: lift a diagonal loop and correct the endpoint
in the connected configuration fiber. Its kernel
already comes from that fiber in the divisor.
This is the same exact comparison proved in F.20,
now with old label $u$ and new label $\epsilon a$.
The omitted labels are still three old points and
one new point; their distinctness and avoidance
conditions give those same connected fibers.
It proves the corresponding discrete version of
(F.21.2) for each of the four rows of (F.21.12)
and both choices of sign.

On full completions, F.19--F.20 identify this precise
divisor image with the inverse image of the full
diagonal subgroup, and its full coset space with
$\widehat F_3$. The actual configuration sections
and characteristic finite quotients of the finitely
generated free kernels used there apply with these
same omitted labels. Thus no completion of an
unproved intersection is being substituted here.
The orbit of a point under the compact group
$L=\widehat{\pi_1(Y)}$ is compact and closed in
$\widehat F_3$. By the discrete transitivity just
proved it contains the dense image of $F_3$.
It is therefore all of $\widehat F_3$.
Transport by $\gamma(\omega)$ gives exactly
(F.21.2) at the original F.19 basepoint.

\emph{The tubular centralizer lift.}
Write $D_i$ for one selected cross divisor in $X$
and $D_i^{\mathrm{reg}}$ for the complement in it
of its intersections with the other seven cross
divisors. This open submanifold is connected:
$D_i$ is connected by the configuration-fiber
description above, and paths can be perturbed
away from its finite union of complex divisors.
The same perturbation makes
$\pi_1(D_i^{\mathrm{reg}})\to\pi_1(D_i)$ onto.
The point $\gamma(d)$ belongs to
$D_i^{\mathrm{reg}}$ by the distinctness of the
sixteen source collision points.

Choose a Hermitian metric on the complex normal
line to $D_i^{\mathrm{reg}}$. A sufficiently small
positive continuous radius in that line gives
a circle bundle lying in the complement of all
eight divisors, namely the ten-point moduli space.
Such a radius is obtained from a locally finite
system of tubular coordinate charts by a positive
partition-of-unity minorant of their admissible
radii. Over any compact loop the resulting
normal circles and their homotopies stay in
these tubular charts.

The image of the fiber generator is central
in the fundamental group of this circle bundle.
Indeed transport around a base loop acts on
$\pi_1(S^1)=\mathbb Z$ by $+1$, because the normal
line is complex oriented. The circle-bundle
homotopy sequence is onto on the base fundamental
group, since the fiber is connected.
Base this construction by the chosen source
whisker and its small terminal normal circle.
Its fiber generator maps to exactly
$\nu_{u,\epsilon}$: F.17.7 gives the simple
holomorphic collision, and hence positive exponent
one. Filling the divisor takes each of these
circle-bundle loops to its base loop in $D_i$.
Consequently the discrete centralizer of this
very meridian maps onto the entire discrete
divisor image in $\pi_1(X)$.

Compose the circle-bundle homomorphism with the
full completion $G$. It extends continuously
from the full completion of the circle-bundle
group. Its compact image centralizes
$\nu_{u,\epsilon}$ and its image in $\mathcal Q$
contains the closure of the whole discrete
divisor image. That closure is exactly
$\mathcal P_{u,\epsilon}$ by the specified
completion comparison above. This proves
(F.21.3). Injectivity of the circle-bundle
completion is not used.

\emph{Normal generation in the actual kernel.}
Let $N_D$ be the closed normal closure of $K^+$
inside $D^+$. Fix one selected meridian $\nu_i$.
For $g\in G$, transitivity (F.21.2) factors its
image as
\[
                 \varpi(g)=\gamma(l)p,
          \qquad l\in L,\quad p\in\mathcal P_i.
\]
Choose $h\in H^\dagger$ over $l$ using the onto
map in F.19.10, and choose
$c\in C_G(\nu_i)$ with $\varpi(c)=p$ using
(F.21.3). Then $d=g c^{-1}h^{-1}\in D^+$ and
\[
                  g=d h c,
       \qquad g\nu_i g^{-1}=d(h\nu_i h^{-1})d^{-1}\in N_D.
                                                               \tag{F.21.27}
\]
Here $\nu_i\in K^+$ and $K^+\triangleleft H^\dagger$
by the exact source row F.19.10, so the middle
conjugate belongs to $K^+$.
F.19.3 states that the closed ambient normal
closure in $G$ of these eight $\nu_i$ is $D^+$.
Equation (F.21.27) puts every such conjugate in
$N_D$, proving $D^+=N_D$.

In the quotient $D^+_{\mathrm{ab}}$, conjugation
by $D^+$ acts trivially. The image of $K^+$ is
a closed subgroup, since $K^+$ is compact.
The normal-generation equality therefore says
that this image is the entire abelianization.
For any continuous finite nilpotent quotient
$q:D^+\twoheadrightarrow P$, put $S=q(K^+)$.
The proved equality says that $S$ normally
generates $P$. Every maximal proper subgroup
$M_0$ of a finite nilpotent group is normal:
choose the first term $Z_j(P)$ of its upper
central series not contained in $M_0$.
Then $Z_{j-1}(P)\subseteq M_0$ and an element
of $Z_j(P)\setminus M_0$ normalizes $M_0$,
because its commutators with $M_0$ lie in
$Z_{j-1}(P)$. The normalizer is strictly
larger than the maximal subgroup and hence
is all of $P$. A proper $S$ would be contained
in such a normal $M_0$; its normal closure
would remain there. Therefore $S=P$.
Finally the image of compact $K^+$ in
$D^+_{\mathrm{nil}}$ is closed and maps onto
every defining finite quotient. It is dense
by the inverse-limit topology and thus equals
$D^+_{\mathrm{nil}}$. This proves all the
stated quotient surjections.

The exact symbolic certificate
\path{code/part_a/cross_divisor_geometry_f21.py}
checks (F.21.11)--(F.21.12), the coordinate
inverse in (F.21.10), the divisor derivative
(F.21.14), the section pullbacks, and all branch
and puncture identities (F.21.20)--(F.21.21).
Its finite checks supplement the bundle,
fundamental-group, and compact-orbit arguments
written here; those arguments establish the
full assertions.
\end{proof}
